\documentclass[10pt]{article}
\usepackage[margin=20mm]{geometry}
\usepackage{amsmath,amssymb,fontspec,unicode-math,xurl,hyperref}
\IfFontExistsTF{DejaVu Serif}{\setmainfont{DejaVu Serif}}{\setmainfont{Latin Modern Roman}}
\setmathfont{latinmodern-math.otf}
\setlength{\parindent}{0pt}\setlength{\parskip}{6pt}\sloppy
\begin{document}
Abstract: The published DALI objective, its minimizer argument, and its MMD concentration argument were checked against the original PDF. An explicit Gaussian-ReLU construction has an attained global minimum but a strictly positive distribution-alignment loss; thus existence of a minimizer does not imply the simultaneous alignment asserted in Proposition 4.5. Independently, the absolute-deviation inequality used to prove Proposition 4.4 is false for a bounded singleton network class. Both defects affect the theoretical justification of the proposed specification objective and support Major impact. The reported benchmark improvements are not refuted by these calculations. Scientific adjudication remains pending.\par
Author: AI-assisted technical review\par
Date:\par
Title: Technical Verification of
Learnware Specification via Dual Alignment\par
\# Identity, version, and inspected scope\par
Wei Chen, Jun-Xiang Mao, Xiaozheng Wang, and Min-Ling Zhang, Learnware Specification via Dual Alignment, ICML 2025, PMLR 267, pp. 8683–8699; local identifier I25-43500 (0131), OpenReview identifier `z3qwOvWbC8`. This review concerns the staged, published 17-page proceedings PDF, not a later arXiv revision. Page references below use one-based PDF pages, which coincide with the printed internal numbers 1–17.\par
The inspection covered the abstract and identity; Sections 3–4, including objectives (2)–(4), Propositions 4.3–4.6 and deployment formulas; Section 5’s experimental definitions and Tables 1–4; Algorithms 1–3; and Appendices B–E. References on pp. 9–10 were not systematically verified. The central calculations below concern Proposition 4.5 and Appendix B.2. Proposition 4.3 was read but its external learning-theory citation was not independently reconstructed. The source manifest reports no identified released repository and an OpenReview metadata access failure; neither absence proves that no additional material exists. No author code, datasets, weights, or training runs were executed.\par
\# Definitions and dependencies\par
With class labels fixed, equation (3) uses the expectation over random embedding networks of the squared difference of empirical class means. Equation (4) minimizes its sum with the smoothed loss discrepancy in equation (2), where \(g(t)=\sqrt{t^2+\alpha}\) and \(\alpha>0\). Assumption 4.1 permits a nonlinear random embedding network and a classifier with the same hidden architecture followed by a linear classification head. Assumption 4.2 assumes that the objective has a minimizer. Proposition 4.5 also initializes the reduced set from training observations. None of these assumptions states that the reduced set can represent every embedded training mean exactly.\par
Appendix B.2 instead introduces a supremum over a function class, with empirical Rademacher complexity defined in equation (20). That supremum is an unsquared integral probability metric, whereas equation (3) is an expected squared feature-mean discrepancy. The distinction matters when transferring a concentration bound to the actual objective.\par
\# Substantive findings\par
Problem 1 (An attained minimum need not give simultaneous alignment).  Assessment: Major. Proposition 4.5, p. 5, equation (8), asserts agreement of model performance and barycenters at the optimized specification; Appendix C, pp. 14–17, culminates in equation (45), asserting convergence of the corresponding expression to zero. The argument uses gradients of the two losses and then infers preservation and alignment. A descent direction and existence of a minimizer do not imply a zero residual.\par
\#\# Construction and verification of assumptions\par
Take \(d=100\), \(K=1\), and a training multiset \(D\) containing 50 copies of \(e_1\) and 50 copies of \(-e_1\). Every true and pseudo-label is the single class. Set \(|R|=1\), so the specification is one vector \(z\in\mathbb R^{100}\), initially sampled as \(e_1\) or \(-e_1\) from \(D\). The embedding is \[\psi_W(x)=(W^\top x)_+,\qquad W\sim N(0,I_{100}).\] This is a nonlinear random ReLU feature with one output. A classifier with the same hidden architecture and a one-logit softmax has class probability one for every input, whatever its linear-head weights. Consequently every categorical cross-entropy is zero. The discriminative term is the constant \(\sqrt{\alpha}\) for all \(z\); taking either the stated supremum or an expectation does not change this. The construction uses the explicitly permitted one-class simplification. It is not an assumption that binary sigmoid loss equals one-class categorical loss.\par
The reduced set is much smaller than the training sample, initialization uses an actual training point, and the stipulated global minimizer will be exhibited. Repeated observations are allowed in an indexed training sample; no distinctness or nondegeneracy assumption appears in the proposition. If distinct points or at least two classes were intended, they would be additional hypotheses, not a consequence of Assumptions 4.1–4.2.\par
\#\# Independent calculation\par
Equation reference: eq:lower. The training feature mean is \(|W_1|/2\). Write \(r=\|z\|_2\). By central symmetry of the Gaussian vector, \[\mathbb E\{|W_1|(W^\top z)_+\}
 =\tfrac12\mathbb E\{|W_1||W^\top z|\}\leq \tfrac r2.\] The inequality is Cauchy–Schwarz, using \(\mathbb E W_1^2=1\) and \(\mathbb E(W^\top z)^2=r^2\). Therefore \[\begin{aligned}
 \mathcal L_{\rm MMD}(D,\{z\})
 &=\mathbb E\bigl[|W_1|/2-(W^\top z)_+\bigr]^2\\
 &\geq \tfrac14+\tfrac12r^2-\tfrac12r
 =\tfrac18+\tfrac12(r-\tfrac12)^2. \end{aligned}\] Equality is attained at \(z=e_1/2\) and \(z=-e_1/2\). Thus Assumption 4.2 is satisfied: the full objective has global minimum \(\sqrt{\alpha}+1/8\). Nevertheless its distribution-alignment component never vanishes. For \(z=e_1/2\), every \(W_1<0\) yields a strictly positive feature-mean residual \(|W_1|/2\), an event of probability \(1/2\). The input barycenter is also different: \(D\) has mean zero and \(R\) has mean \(e_1/2\).\par
This refutes exact simultaneous alignment, even at a global optimum. It does not depend on failure of Adam, inadequate iteration count, or a bad local minimum. The independently written `independent\_check.py` verifies the completed-square identity; `independent\_check.json` records the parameters and exact minimum. The probabilistic identities and inequality above are analytic, not inferred from simulation.\par
\#\# Counterevidence and consequence\par
An approximate interpretation of the method remains coherent: minimizing a nonnegative loss may improve alignment without making it exact. Figure 3 is evidence about its plotted experiment, not a proof of universal equality. The construction does not challenge a theorem restricted to exactly representable distributions, but such a restriction is absent here. The literal equation (8) also combines signed residuals; our classifier residual is zero, so cancellation with a classification term cannot rescue it.\par
Proposition 4.5 is the theoretical minimization claim for the central dual-alignment objective and is then used explicitly in Proposition 4.6’s privacy argument. This core dependency supports Major impact. An appropriate narrower claim would acknowledge an approximation residual, but no general convergence or privacy repair is established in this review.\par
Problem 2 (The MMD proof drops a necessary absolute deviation).  Assessment: Major. Appendix B.2, p. 13, equation (24), upper-bounds the absolute difference of the population and empirical suprema by a supremum of their signed differences. This step is not valid for the general network classes stated there. It feeds equations (25)–(28) and Proposition 4.4, p. 5, equation (7).\par
\#\# Complete finite counterexample\par
Let \(\mathcal F=\{f\}\) with \(f(x)=x_+\) on inputs \(\{0,1\}\). This is a ReLU network of depth one, weight norm one, bounded input norm one, with a 1-Lipschitz positively homogeneous activation. A deterministic parameter distribution is a permissible special case; no symmetry or richness requirement is stated for the class in equation (20).\par
Choose \(p=\operatorname{Bernoulli}(3/4)\) and \(q=\operatorname{Bernoulli}(1/4)\), independently sample \(D=\{1\}\) and \(R=\{0\}\), and condition on this event, whose probability is \(9/16\). The two distributions are different, as required immediately before equation (24). The population supremum is \(1/2\) and the empirical supremum is \(1\). The left side of (24) is consequently \(1/2\), but the proposed right side \(\Delta(p,q,D,R)\) is \(-1/2\). The asserted inequality would require \(1/2\leq-1/2\).\par
The original PDF image of p. 13 was inspected: this is the inequality as typeset, not a lost absolute-value symbol in text extraction. A valid general elementary inequality is \[\left|\sup_{f\in\mathcal F}a_f-\sup_{f\in\mathcal F}b_f\right|
 \leq \sup_{f\in\mathcal F}|a_f-b_f|,\] assuming finite suprema. This repairs only that step; it does not verify the subsequent concentration argument or convert the unsquared supremum into equation (3).\par
\#\# Complexity interpretation and further check\par
On the literal empirical-complexity definition (20), a singleton class has \[\widehat{\mathcal R}_S(\{f\})
 =\mathbb E_\sigma\,|S|^{-1}\sum_i\sigma_i f(x_i)=0.\] Under \(p=q=\operatorname{Bernoulli}(1/2)\), take \(|D|=100\), \(|R|=1\), all observations independent. The two empirical means differ with probability \(1-2^{-100}\), so their squared difference is positive with that probability. The right side of (7), using those actual complexities, is zero. For \(\delta=1/4\), the asserted high-probability bound fails. The sample sizes also respect substantial reduction.\par
There is an alternative reading: the symbols might stand for the positive norm-based upper bounds in (21), rather than the complexities defined in (20). That reading avoids this particular zero-complexity counterexample, but does not validate the signed inequality (24), establish the product-form concentration term, or justify applying an independent-sampling theorem to a specification optimized from \(D\). The supported assessment is an invalid central derivation, plus falsity under the literal complexity definition; it is not a claim that every suitably revised bound must be false.\par
\#\# Impact\par
Proposition 4.4 supplies the distribution-alignment half of the claimed theoretical support, expressly combined with Proposition 4.3 below equation (7). Its proof therefore has a direct core dependency. Rebuilding the correct function class, bounded-difference constants, square/norm conversion, and adaptive-sampling treatment requires a fresh theorem. Published accuracy tables are not recalculated or invalidated by this finite check.\par
\# Overall outcome and unresolved scope\par
Both Problems concern the mathematical support for DALI’s central objective. Empirical comparisons and the usefulness of optimizing a compromise objective may survive. Appendix D’s privacy discussion depends on Proposition 4.5 and gives no formal threat model that would turn small specification size into a general privacy guarantee. Its membership-inference experiment in Section 5.4 is an empirical result with its own scope; no new privacy attack was run here.\par
No later-version comparison or global erratum search was performed; the inspected official metadata identifies the published proceedings version. There is no basis in this review to label these first discoveries or known author-corrected issues. No proposed repaired DALI theorem or implementation has been verified.\par
Item | Issue assessment | Repair verification | Adoption / remaining work\par
Proposition 4.5 | Major; explicit positive optimum | No general repair | Unknown; quantify attainable residual\par
MMD derivation | Major; false inequality | One elementary inequality only | Unknown; rebuild bound and sampling assumptions\par
Experiments | Not reproduced | Not applicable | Data, code, and fresh runs would be needed\par
Scientific decision | Pending | Separate from issue judgment | No human review supplied\par
1 W. Chen, J.-X. Mao, X. Wang, and M.-L. Zhang, Learnware Specification via Dual Alignment, Proceedings of ICML, PMLR 267 (2025), 8683–8699. Published proceedings PDF; https://proceedings.mlr.press/v267/chen25as.html.\par
\# Evidence, reproducibility, and roles\par
The immutable original has SHA-256\par
> .\par
The accompanying numeric verification directory contains source-integrity facts, the extracted text, original-page images 5, 13, and 17, and the independent Python check with its actual output. The script uses SymPy only for an exact polynomial identity and rational arithmetic for the finite example. It does not import or execute author code. The JSON ledger distinguishes reading from analytic verification, original-PDF visual checks, and any report compilation/visual status. Full command and hash bindings are in `provenance.json`; unavailable operations are explicitly recorded there.\par
This is an AI initial scientific review in the assigned conversation. Human review has not been supplied; `human\_reviews` remains empty and adjudication is pending. No author adoption was requested.\par
\end{document}
